\begin{papercard}{arksey2005}
\cf{Stages} \stagetag{S} (what a scoping study is and why use one), \stagetag{A} (question, search sources, eligibility and selection), \stagetag{D} (``charting'' the data), \stagetag{R} (collating and reporting without synthesis). Arksey and O'Malley, ``Scoping studies: towards a methodological framework'' (\emph{International Journal of Social Research Methodology} 8(1):19--32, 2005; here the 31-page White Rose author version). The first widely cited framework for scoping studies: five stages plus an optional consultation, illustrated with the authors' own scoping study of services for carers of people with mental health problems. It is a conduct framework, written before reporting guidance existed; PRISMA-ScR (Tricco et al.) is the reporting counterpart.

\cf{Position in the literature-review family} The authors note a ``plethora of terminology'' (systematic review, meta-analysis, rapid review, narrative review, structured review) with no consistent definitions, so labels are used loosely; a review called ``systematic'' without the same protection against bias should be called a literature review. A scoping study is another type of review: a systematic review focuses on a well-defined question with study designs identifiable in advance, a scoping study addresses broader topics where many designs apply \pc{arksey2005}{7}{review might typically focus on a well defined question where appropriate study designs can be identified in advance whilst a scoping study tends to address broader topics}; they argue that ``there is no single `ideal type' of literature review'' but a set of tools to be used appropriately \pc{arksey2005}{7}{there is no single ideal type of literature review}. The abstract describes scoping as ``an approach to reviewing the literature which to date has received little attention'' \pc{arksey2005}{5}{an approach to reviewing the literature which to date has received little attention}.

\cf{Definition and four purposes} They quote Mays et al.: scoping studies ``aim to map rapidly the key concepts underpinning a research area and the main sources and types of evidence available'' \pc{arksey2005}{8}{aim to map rapidly the key concepts underpinning a research area and the main sources and types of evidence available}. Breadth (coverage) is distinguished from depth (information extracted per study). Four reasons to do one: (1) \emph{examine the extent, range and nature of research activity} \pc{arksey2005}{9}{To examine the extent, range and nature of research activity}; (2) \emph{determine the value of a full systematic review} (does literature exist, has a review been done, what would it cost) \pc{arksey2005}{9}{To determine the value of undertaking a full systematic review}; (3) \emph{summarise and disseminate findings} to people lacking time; (4) \emph{identify research gaps} \pc{arksey2005}{10}{To identify research gaps in the existing literature}. Types 1--2 are steps towards a systematic review, types 3--4 methods in their own right; the paper concentrates on type 4. Gaps found this way are gaps of \emph{absence}: ``quality assessment does not form part of the scoping study remit'', so poor-quality research areas are not flagged \pc{arksey2005}{10}{quality assessment does not form part of the scoping study remit}.

\cf{Principles} The framework rests on the systematic-review demand for rigour and transparency: the process ``should be documented in sufficient detail to enable the study to be replicated by others'' \pc{arksey2005}{11}{The process should be documented in sufficient detail to enable the study to be replicated by others}. Because the aim is to identify all relevant literature regardless of design, researchers may not set strict limits on search terms or selection at the outset, and as familiarity grows they redefine terms \pc{arksey2005}{11}{not wish to place strict limitations on search terms, identification of relevant studies, or study selection at the outset}. The process ``is not linear but iterative'' and reflexive, repeating steps where needed \pc{arksey2005}{11}{The process is not linear but iterative, requiring researchers to engage with each stage in a reflexive way}.

\cf{Stage 1: identifying the research question} Choose the facets that matter (population, intervention, outcomes) because the question guides the search. Their question asked what is known about effectiveness and cost-effectiveness of services supporting carers; the term ``services to support carers'' was ambiguous and ``mental health problems'' needed a decision on which illnesses to include. Wide definitions reduce the chance of missing papers but generate unmanageable numbers; recommendation: ``maintain a wide approach in order to generate breadth of coverage'' and tighten parameters once the volume is known \pc{arksey2005}{13}{Our recommendation would be to maintain a wide approach in order to generate breadth of coverage}.

\cf{Stage 2: identifying relevant studies} The point is to be as comprehensive as possible \pc{arksey2005}{13}{the whole point of scoping the field is to be as comprehensive}, using electronic databases, reference lists, hand-searching of key journals, and networks, organisations and conferences. Practical limits are set at the outset: only studies published January 1985 to October 2001 (policy changes; budget) and no foreign-language material because of translation cost; ``potentially relevant papers could have been missed'' \pc{arksey2005}{14}{potentially relevant papers could have been missed}. For databases a CRD information officer, a qualified librarian, devised an initial strategy that was refined after early results, first run on MEDLINE and then converted for the other databases \pc{arksey2005}{14}{She then devised an initial search strategy, which was later refined in the light of early results}; researchers may lack such skills \pc{arksey2005}{14}{Researchers may not have the skills necessary for designing and executing sensitive search strategies that qualified librarians have}. Searches covered 12 CD-ROM databases and four internet ones, with huge variation (four gave under ten hits each; MEDLINE 1,565 and EMBASE 1,589) \pc{arksey2005}{15}{searches were made on 12 databases available from CD-ROMs, and four via the Internet}; it is unknown which database yielded the 204 included articles \pc{arksey2005}{15}{It is not known how effective the different databases were in generating the 204 articles}. Reference-list checking reached ``a saturation point'' where nothing new appeared \pc{arksey2005}{15}{although a saturation point was reached where no new ones were being identified}; citation searching was not used. Hand-searching is needed because databases are incomplete and ``they all tend to have a Western and particularly US bias'' \pc{arksey2005}{15}{they all tend to have a Western and particularly US bias}. Four journals were hand-searched; organisations (Carers UK, Mental Health Foundation, King's Fund and others) were contacted. Records were managed in EndNote; the information officer logged each database, years covered and search date \pc{arksey2005}{16}{The Information Officer recorded each database searched, the years it covered, and the date it was searched}, important for updating. Total yield: 3,867 references, 112 found during the study, 3,755 from databases \pc{arksey2005}{17}{generated a total of 3,867 references, some 112 of which were identified as the study progressed}.

\cf{Stage 3: study selection} The search captured many irrelevant studies because breadth was sought and country terminology differed. Systematic reviews fix inclusion and exclusion criteria beforehand; here ``criteria were devised post hoc, based on increasing familiarity with the literature'' \pc{arksey2005}{17}{criteria were devised post hoc, based on increasing familiarity with the}, relating to type of study, type of intervention, care-recipient group and carer group. Two reviewers applied the criteria to all citations \pc{arksey2005}{18}{Two reviewers then applied the inclusion and exclusion criteria to all the citations}; full articles were obtained for ``best fit'' and unclear cases; a deadline cut off further additions, with unreviewed items best listed in an appendix \pc{arksey2005}{18}{A deadline was set, after which it was agreed that we would not include any more studies}. Of 3,867 references 453 were ordered through inter-library loan and 204 were included after full reading \pc{arksey2005}{18}{Out of our original 3,867 references, 453 were ordered through inter-library loans}. No agreement statistic between the two reviewers is reported.

\cf{Stage 4: charting the data} ``Charting'' is borrowed from Ritchie and Spencer: a technique for synthesising and interpreting qualitative data ``by sifting, charting and sorting material according to key issues and themes'' \pc{arksey2005}{18}{describes a technique for synthesising and interpreting qualitative data by sifting, charting and sorting material according to key issues and themes}; in a systematic review the same step is called data extraction \pc{arksey2005}{19}{this process would be called data extraction and, in the case of meta-analysis, might involve specific statistical techniques}. The approach is a ``descriptive-analytical'' narrative method: a common analytical framework is applied to all reports, collecting standard information on each \pc{arksey2005}{19}{involves applying a common analytical framework to all the primary research reports and collecting standard information on each study}. A short summary per study alone does not help decision makers. Data went into an Excel ``data charting form'' \pc{arksey2005}{19}{entered onto a data charting form using the database programme Excel}: authors, year, location; intervention type, comparator, duration; study populations; aims; methodology; outcome measures; important results, plus economic data where present. A uniform approach was sought for all 204 studies, but it was often impossible to extract everything because reports omitted material \pc{arksey2005}{20}{it was often impossible to extract all the information required where research reports failed to include relevant material}.

\cf{Stage 5: collating, summarising and reporting} Unlike systematic reviews, where a small percentage of studies reach the report and the rest remain hidden, the scoping study ``seeks to present an overview of all material reviewed'' \pc{arksey2005}{20}{the scoping study seeks to present an overview of all material reviewed}. It does not synthesise or aggregate findings or state the ``weight'' of evidence \pc{arksey2005}{21}{the scoping study does not seek to synthesise evidence or to aggregate findings from different studies}, because quality is not assessed and robustness cannot be judged. Reporting took two forms: basic numerical analysis of extent, nature and distribution (geography, care-recipient group, intervention range, methods, measures) \pc{arksey2005}{21}{attention was given to basic numerical analysis of the extent, nature and distribution of the studies included in the review}, and a thematic organisation into eleven intervention types, which became ``the primary unit of analysis'' despite overlap and inconsistent author definitions \pc{arksey2005}{21}{The intervention type became the primary unit of analysis}. Choosing that frame made findings ``subsume theoretical or conceptual positions adopted by authors'' (an alternative would have been competing theories), so the reporting strategy must make potential bias visible \pc{arksey2005}{22}{our findings tended to subsume theoretical or conceptual positions adopted by authors}. A fixed template per group began with a table and then nine headings: interventions, sample sizes, participants, research methods, outcomes, evidence on effectiveness, economic aspects, UK studies, gaps \pc{arksey2005}{22}{commentary written under the following nine headings: interventions; sample sizes; participants; research methods; outcomes}.

\cf{Optional stage: consultation} Practitioners and users can improve the work; the team consulted national bodies, local managers and practitioners, and ``key informant'' carers \pc{arksey2005}{23}{involved three groups of stakeholders: representatives from national statutory and}. They supplied extra references and a broader view of effectiveness (five dimensions including benefit to the family and long-term societal outcomes), which led the team to question reliance on standard outcome measures \pc{arksey2005}{24}{Although this element of our approach to a scoping study may be considered an optional extra, the consultation exercise did indeed provide added value}.

\cf{Resources, strengths and weaknesses} The review used three full-time-equivalent staff for six months plus an information officer \pc{arksey2005}{24}{employed three full time equivalent staff members for six months}; a scoping study ``should not be seen as a cheap alternative'' \pc{arksey2005}{25}{scoping study should not be seen as a cheap alternative to the systematic review}. The authors stress there is ``no definitive procedure'' \pc{arksey2005}{25}{there is no definitive procedure for scoping the literature}. Strength: a transparent way to map the volume, nature and characteristics of research and identify gaps. Limits they name: no formal appraisal of evidence quality \pc{arksey2005}{26}{They do not, for example, appraise the quality of evidence in the primary research reports in any formal sense}, a trade-off between breadth and depth, no synthesis, so the result is ``a narrative or descriptive account'' \pc{arksey2005}{26}{Consequently, scoping studies provide a narrative or descriptive account of available research}. It is not an easy option: high analytic skill is needed to frame large numbers of studies \pc{arksey2005}{27}{Conducting a scoping study requires reviewers to have high degrees of analytic skill}, and there is no clear means of synthesising across designs \pc{arksey2005}{27}{the scoping study does not offer any clear means of synthesising findings from different kinds of study design}.

\cf{Limits (card)} A single worked example (UK carer services, searches to 2001) described by its own authors; the framework is a description of what they did, not a validated protocol; post hoc criteria and a single deadline-based stop are acknowledged as pragmatic; reviewer agreement is not measured; the framework predates PRISMA-ScR, PRISMA-S and modern LLM-assisted screening; it assumes bibliographic databases and says nothing about app stores, web documents or privacy policies as ``sources''; the consultation step is defined narrowly for health and social care.

\cf{Use for this thesis} (1) A citable basis for calling the literature and app part a \emph{mapping} exercise rather than a systematic review: the thesis can state that it describes extent, range and nature of how harms are communicated, does not weigh evidence or rate app quality, and flags gaps of absence only. (2) Stage 1 advice (start wide, define ambiguous terms such as ``harm'', ``health app'', ``privacy communication'' at the outset, narrow after seeing the volume) matches the scope work in stage S. (3) Stage 3 admits post hoc criteria; for the thesis it is cleaner to fix eligibility before drawing the app sample and to log every later change in the Decision Log, using Arksey and O'Malley only as precedent. (4) ``Charting'' with a uniform form is the model for the app data sheet (store, version, date, category, developer, permissions, label content, policy URL); missing fields should be recorded as missing, as they note for incomplete reports. (5) A fixed reporting template per group (here: per harm type or per stage) supports comparison across chapters and makes reporting bias visible. (6) The consultation stage suggests, if time allows, a short expert or supervisor check of the harm tree; mark it optional. (7) Pair with PRISMA-ScR for what to report and PRISMA-S for the search log.
\end{papercard}
